\documentclass[12pt,a4paper]{report}

% =========================================================================
% CONFIGURATION: ENGLISH LANGUAGE & GENUINE TIMES NEW ROMAN FONT SYSTEM
% =========================================================================
\usepackage[utf8]{inputenc}
\usepackage[T5,T1]{fontenc}
\usepackage[vietnamese,english]{babel}
\usepackage{mathptmx}

% =========================================================================
% PAGE GEOMETRY AND DOCUMENT MARGINS
% =========================================================================
\usepackage{geometry}
\geometry{
	left=3cm,
	right=2cm,
	top=2.5cm,
	bottom=2.5cm
}

% =========================================================================
% GRAPHICS, UTILITIES, AND GRAPH ENVIRONMENT
% =========================================================================
\usepackage{graphicx}
\usepackage{grffile}
\usepackage{tikz}
\usepackage{float}
\usepackage{setspace}
\usepackage{indentfirst}
\usepackage{titlesec}
\usepackage{booktabs}
\usepackage{longtable}
\usepackage{xcolor}
\usepackage{colortbl}
\usepackage{listings}
\usepackage{enumitem}
\usepackage{amsfonts,amsmath,amssymb}
\usepackage{tcolorbox}
\usepackage{hyperref}
\usepackage{array}
\usepackage{multirow}
\usepackage{tabularx}
\usepackage{fancyhdr}

\usetikzlibrary{calc, shapes.geometric, arrows.meta, positioning, backgrounds, fit, automata, shapes.multipart, decorations.pathreplacing}

\setlength{\parindent}{1.27cm}

% =========================================================================
% COLOR DEFINITIONS
% =========================================================================
\definecolor{mainblue}{RGB}{26, 82, 118}
\definecolor{sectionblue}{RGB}{33, 97, 140}
\definecolor{titlecolor}{RGB}{180, 0, 0}
\definecolor{primaryblue}{RGB}{26, 82, 118}
\definecolor{lightgray}{RGB}{240, 240, 240}
\definecolor{accentgreen}{RGB}{39, 174, 96}
\definecolor{accentorange}{RGB}{230, 126, 34}

% =========================================================================
% CODES AND SCRIPT LISTINGS CONFIGURATION
% =========================================================================
\definecolor{bluekeywords}{rgb}{0.13,0.13,1}
\definecolor{greencomments}{rgb}{0,0.5,0}
\definecolor{redstrings}{rgb}{0.9,0,0}
\definecolor{graynumbers}{rgb}{0.5,0.5,0.5}

\lstset{
	backgroundcolor=\color{gray!5},
	commentstyle=\color{greencomments},
	keywordstyle=\color{bluekeywords}\bfseries,
	stringstyle=\color{redstrings},
	numberstyle=\tiny\color{graynumbers},
	basicstyle=\ttfamily\footnotesize,
	breakatwhitespace=false,
	breaklines=true,
	captionpos=b,
	keepspaces=true,
	numbers=left,
	numbersep=5pt,
	showspaces=false,
	showstringspaces=false,
	showtabs=false,
	tabsize=2,
	extendedchars=true,
	language=SQL
}

% =========================================================================
% HEADINGS AND TITLES FORMATTING
% =========================================================================
\renewcommand{\thechapter}{\Roman{chapter}}

\titleformat{\chapter}[display]
{\normalfont\bfseries\Large\color{mainblue}}
{\chaptertitlename\ \thechapter}{14pt}{\Large}

\titleformat{\section}
{\normalfont\bfseries\fontsize{14}{18}\selectfont\color{black}}
{\arabic{chapter}.\arabic{section}}{1em}{}

\titleformat{\subsection}
{\normalfont\itshape\fontsize{13}{17}\selectfont\color{black}}
{\arabic{chapter}.\arabic{section}.\arabic{subsection}}{1em}{}

\titleformat{\subsubsection}
{\normalfont\fontsize{12}{16}\selectfont\color{sectionblue}}
{\arabic{chapter}.\arabic{section}.\arabic{subsection}.\arabic{subsubsection}}{1em}{}

% =========================================================================
% HYPERREF SETUP
% =========================================================================
\hypersetup{
	colorlinks=true,
	linkcolor=mainblue,
	urlcolor=mainblue,
	citecolor=mainblue,
	pdftitle={Personalized Career Orientation & Learning Roadmap Platform},
	pdfauthor={Group 2 - UTH}
}

% =========================================================================
% DOCUMENT INITIALIZATION
% =========================================================================
\begin{document}
	
	% =========================================================================
	% PAGE 1: OFFICIAL COVER PAGE
	% =========================================================================
	\begin{titlepage}
		\begin{center}
			
			\begin{tikzpicture}[remember picture,overlay]
				\draw[line width=2pt, draw=mainblue]
				($(current page.north west)+(1.5cm,-1.5cm)$)
				rectangle
				($(current page.south east)+(-1.5cm,1.5cm)$);
				
				\draw[line width=0.8pt, draw=mainblue]
				($(current page.north west)+(1.7cm,-1.7cm)$)
				rectangle
				($(current page.south east)+(-1.7cm,1.7cm)$);
			\end{tikzpicture}
			
			\vspace*{0.5cm}
			
			{\Large\textbf{UNIVERSITY OF TRANSPORT HO CHI MINH CITY}}\\
			{\large\textbf{INSTITUTE OF HIGH QUALITY TRAINING}}
			
			\vspace{1.2cm}
			
			\begin{center}
				\makeatletter
				\IfFileExists{"logo UTH.jpg"}{\includegraphics[width=9cm]{"logo UTH.jpg"}}{\framebox[9cm]{\rule{0pt}{3cm}\textit{[UTH Logo]}\rule{0pt}{3cm}}}
				\makeatother
			\end{center}
			
			\vspace{0.8cm}
			
			{\Large\textbf{FINAL TERM REPORT}}\\[0.3cm]
			{\Large\textbf{SOFTWARE ENGINEERING COURSE}}
			
			\vspace{1.0cm}
			
			{\Large\textbf{PROJECT TOPIC}}
			
			\vspace{0.5cm}
			
			{\fontsize{18}{24}\selectfont\bfseries\color{titlecolor}
				PERSONALIZED CAREER ORIENTATION\\
				\medskip
				\& LEARNING ROADMAP PLATFORM\\
				\medskip
				FOR SOFTWARE ENGINEERING STUDENTS\par}
			
			\vspace{0.8cm}
			
			{\large\textbf{Group 2}}
			
			\vfill
			
			\begin{tabular}{ll}
				\textbf{Instructor:} & Nguyen Van Chien \\[0.5cm]
				\textbf{Students:}   & Vo Trong Nguyen \\
				& Nguyen Van Bao \\
				& Vo Thi Kiem Dung \\
				& Vo Bao Phuc \\
				& Tran Thi Bich Diem \\
				& Phan Thi Duy Chung \\
				& Vu Le Hoang Nhat \\
			\end{tabular}
			
			\vfill
			
			{\large Ho Chi Minh City, June 2026}
			
		\end{center}
	\end{titlepage}
	
	% =========================================================================
	% PAGE 2: SRS TITLE PAGE
	% =========================================================================
	\begin{titlepage}
		\centering
		\vspace*{1.5cm}
		
		{\Huge\bfseries\color{primaryblue}
			PERSONALIZED CAREER ORIENTATION\\[0.4em]
			\& LEARNING ROADMAP PLATFORM\\[0.3em]
			FOR SOFTWARE ENGINEERING STUDENTS
			\par}
		
		\vspace{1cm}
		\rule{\textwidth}{1.5pt}
		
		\vspace{0.8cm}
		{\large\bfseries\textit{A Software Requirements Specification \& System Design Document}}
		
		\vspace{0.6cm}
		{\normalsize Based on methodologies from:}\\[0.3em]
		{\large\textbf{Software Engineering: A Practitioner's Approach}\\
			Roger S. Pressman, Ph.D. --- 7th Edition}
		
		\vspace{1.2cm}
		\begin{tcolorbox}[colback=lightgray, colframe=primaryblue, width=0.85\textwidth]
			\centering
			\textbf{Project Overview}\\[0.5em]
			This document specifies the requirements, architecture, design models, and\\
			quality assurance strategy for a web-based platform that delivers personalized\\
			career orientation and adaptive learning roadmaps to undergraduate Software\\
			Engineering students, bridging the gap between academic curricula and\\
			real-world industry expectations.
		\end{tcolorbox}
		
		\vfill
		
		{\normalsize
			\begin{tabular}{ll}
				\textbf{Document Type:}  & Software Requirements Specification (SRS) \\[3pt]
				\textbf{Version:}        & 1.0 \\[3pt]
				\textbf{Process Model:}  & Agile/Incremental (Scrum-based) \\[3pt]
				\textbf{Methodology:}    & Pressman's Software Engineering Framework \\[3pt]
				\textbf{Date:}           & \today \\
			\end{tabular}
		}
		
		\vspace{1.5cm}
		\rule{\textwidth}{0.5pt}
	\end{titlepage}
	
	% =========================================================================
	% TABLE OF CONTENTS
	% =========================================================================
	\tableofcontents
	\newpage
	
	% =========================================================================
	% CONTENT SPACING ENVIRONMENT
	% =========================================================================
	\begin{spacing}{1.5}
		
		% =====================================
		% CHAPTER I: PROJECT OVERVIEW
		% =====================================
		\chapter{Project Overview}
		
		\section{Context \& Comprehensive Problem Statements}
		
		The Software Engineering (SE) landscape in 2026 is rapidly evolving, dominated by innovations in Artificial Intelligence, Cloud Native Architectures, DevSecOps pipelines, and Large Language Model (LLM) integrations. Consequently, undergraduate students face an immense pedagogical paradox: academic institutions traditionally educate them to become generalists who grasp foundational computer science theories, yet the commercial labor market aggressively demands highly specialized engineering capabilities upon graduation.
		
		Based on industry surveys and labor market analysis conducted across major Vietnamese technology hubs (Ho Chi Minh City, Hanoi, Da Nang) in Q1--Q2 2026, the gap between academic preparation and employer expectations has widened significantly. Entry-level software engineering positions now routinely require proficiency in cloud-native frameworks (Kubernetes, AWS Lambda), containerization strategies, CI/CD pipeline tooling, and at minimum one specialized domain track such as Machine Learning Operations (MLOps), Backend API Architecture, or Mobile Platform Engineering.
		
		The three core problem clusters motivating this platform are:
		
		\begin{enumerate}
			\item \textbf{Industry-Curriculum Skill Gaps:} Higher education curricula inherently suffer from administrative and structural latencies. While technology stacks iterate within months, university course frameworks take years to adjust through committee reviews, faculty re-skilling, and syllabus approval processes. According to the Stack Overflow Developer Survey 2025, over 67\% of junior developers report that their formal education did not adequately prepare them for their first professional role.
			
			\item \textbf{Information Overload and Choice Paralysis:} The modern internet offers an overwhelming abundance of open-source learning paths, MOOC platforms, YouTube tutorials, and technical documentation. This unstructured wave of materials induces high cognitive fatigue, causing students to either follow generic roadmaps that do not align with their personal strengths, or to abandon structured learning entirely in favor of ad-hoc tutorial consumption.
			
			\item \textbf{Fragmented and Disconnected Portfolios:} The majority of software engineering undergraduates develop minor, uncoordinated applications solely to satisfy immediate course grading criteria. These portfolio artifacts fail to communicate a coherent professional narrative to potential employers. Without strategic guidance, students cannot align their project work with the technical storytelling demanded by competitive recruitment processes at top-tier technology companies.
		\end{enumerate}
		
		\section{System Actors}
		
		The platform identifies four primary stakeholder roles, each interacting with the system across distinct functional domains:
		
		\begin{longtable}{|p{1.5cm}|p{3.5cm}|p{9.5cm}|}
			\hline
			\rowcolor{gray!20}\textbf{No.} & \textbf{System Actor} & \textbf{Functional Role Description} \\ \hline
			1 & Administrator & Manages platform configurations, user credentials, technology skill trees, centralized course catalogs, job market data pipelines, and aggregates global analytics reports. Has full CRUD privileges across all system entities. \\ \hline
			2 & Student & Primary platform consumer. Conducts cognitive skill assessments, receives specialized career orientation paths, tracks personalized learning roadmaps, interacts with the AI Virtual Mentor, manages E-Portfolio, and explores live job market intelligence data. \\ \hline
			3 & Career Advisor & Leverages student performance reports and AI-generated analytical data to execute targeted guidance sessions. Monitors individual academic progress, provides supplementary mentoring notes, and reviews AI-generated recommendations for quality assurance. \\ \hline
			4 & Recruiter & Publishes institutional technology job advertisements, manages job vacancies, applies advanced filtering to source candidates matching target skill profiles, and reviews student E-Portfolio submissions. \\ \hline
		\end{longtable}
		
		\section{System Objectives and Core Innovation}
		
		The overarching goal of this capstone project is to architect and implement an intelligent software engineering ecosystem called the \textbf{Personalized Career Orientation \& Learning Roadmap Platform (PCOLRP)}. The system pursues five strategic objectives:
		
		\begin{enumerate}
			\item \textbf{Personalization at Scale:} Deliver individualized career roadmaps tailored to each student's unique academic transcript, self-assessed competency profile, and stated professional aspirations --- not generic templates.
			
			\item \textbf{Real-Time Market Alignment:} Continuously scrape and analyze live job market data from major IT job portals (LinkedIn, TopCV, ITviec) to ensure that recommended learning paths align with current employer demand signals, not year-old curriculum benchmarks.
			
			\item \textbf{AI-Augmented Mentoring:} Integrate Large Language Model APIs (GPT-4o, Gemini 1.5 Pro) to provide on-demand, context-aware career counseling that supplements limited human advisor bandwidth at scale.
			
			\item \textbf{Portfolio Coherence:} Guide students to build portfolios that tell a deliberate professional story by connecting GitHub repositories, skill acquisitions, and target career roles into a unified, employer-facing narrative.
			
			\item \textbf{Academic-Industry Bridge:} Create a feedback loop between academic institutions and industry trends, enabling career advisors and curriculum designers to identify systemic skill gaps at the cohort level.
		\end{enumerate}
		
		\subsection{Scope Boundaries}
		
		The platform is scoped to serve undergraduate Software Engineering students enrolled at the University of Transport Ho Chi Minh City (UTH) as the primary pilot cohort, with architectural provisions for multi-institution scaling in subsequent development phases. The system explicitly does \textbf{not} replace human academic advisors but augments their capacity through AI-assisted data analytics and personalized recommendation engines.
		
		\section{Background and Theoretical Foundations}
		
		\subsection{Pressman's Software Engineering Framework}
		
		This project applies the structured software engineering methodology outlined by Roger S. Pressman in \textit{Software Engineering: A Practitioner's Approach} (7th Edition). Pressman's process framework prescribes five core workflow activities: Communication, Planning, Modeling, Construction, and Deployment. Each chapter of this document maps directly to these activities, ensuring rigorous traceability from initial problem framing through to implementation and quality validation.
		
		\subsection{Agile/Scrum Development Model}
		
		Given the iterative nature of AI model integration and the need for continuous user feedback during roadmap generation calibration, the project adopts an Agile/Incremental process model organized into two-week Sprints. The Scrum framework governs Sprint Planning, Daily Stand-ups, Sprint Review, and Retrospective ceremonies throughout the 16-week development lifecycle.
		
		\subsection{Related Work and Technology Landscape}
		
		Existing platforms such as Coursera, LinkedIn Learning, and roadmap.sh provide generalized learning paths but lack the personalization depth required for individual student profiles. Academic Career Management Systems (CMS) deployed at large research universities typically address administrative tracking but provide no AI-driven recommendation or real-time market intelligence. This platform addresses the intersection of personalized education technology (EdTech) and career development intelligence that remains underserved in the Vietnamese higher education ecosystem.
		
		% =====================================
		% CHAPTER II: SYSTEM REQUIREMENTS SPECIFICATION (SRS)
		% =====================================
		\chapter{System Requirements Specification (SRS)}
		
		\section{Introduction to the SRS Document}
		
		This Software Requirements Specification (SRS) document is structured in accordance with the IEEE 830-1998 standard and Pressman's requirements engineering framework. It serves as the contractual baseline between stakeholders and the development team, defining what the system must do, the constraints under which it must operate, and the quality attributes it must satisfy.
		
		\subsection{Purpose and Intended Audience}
		
		This SRS targets four distinct reader categories: the development team (for implementation guidance), the project supervisor (for scope validation), future maintenance engineers (for system evolution), and administrative stakeholders (for resource allocation and risk assessment).
		
		\subsection{Definitions, Acronyms, and Abbreviations}
		
		\begin{longtable}{|p{3.5cm}|p{11cm}|}
			\hline
			\rowcolor{gray!20}\textbf{Term / Acronym} & \textbf{Definition} \\ \hline
			SRS & Software Requirements Specification \\ \hline
			FR & Functional Requirement \\ \hline
			NFR & Non-Functional Requirement \\ \hline
			UC & Use Case \\ \hline
			LLM & Large Language Model (e.g., GPT-4o, Gemini 1.5 Pro) \\ \hline
			RBAC & Role-Based Access Control \\ \hline
			JWT & JSON Web Token (authentication mechanism) \\ \hline
			API & Application Programming Interface \\ \hline
			GPA & Grade Point Average \\ \hline
			E-Portfolio & Electronic Portfolio (shareable online project showcase) \\ \hline
			NLP & Natural Language Processing \\ \hline
			CI/CD & Continuous Integration / Continuous Deployment \\ \hline
			PCOLRP & Personalized Career Orientation \& Learning Roadmap Platform \\ \hline
		\end{longtable}
		
		\section{Functional Requirements (FR)}
		
		The functional requirements are organized by the six major subsystem modules identified during stakeholder analysis and use case decomposition workshops.
		
		\subsection{Module FR-1: AI Virtual Mentor}
		
		\begin{longtable}{|p{2cm}|p{4cm}|p{9cm}|}
			\hline
			\rowcolor{mainblue!15}\textbf{FR ID} & \textbf{Requirement Name} & \textbf{Detailed Description} \\ \hline
			FR1.1 & Natural Language Chat Interface & The system must provide a persistent chat interface allowing students to submit career-related queries in natural language (English or Vietnamese). \\ \hline
			FR1.2 & LLM API Integration & The system must integrate with LLM APIs (GPT-4o or Gemini 1.5 Pro) to process student queries and return contextually relevant technical career guidance. \\ \hline
			FR1.3 & Transcript-Aware Personalization & The system must allow students to upload academic transcripts (PDF format) and use extracted GPA and course completion data to contextualize AI responses. \\ \hline
			FR1.4 & GitHub Profile Analysis & The system must retrieve and parse public GitHub profile data (repositories, commit frequency, languages used) to personalize mentoring outputs. \\ \hline
			FR1.5 & Conversation History Persistence & The system must store and display full conversation history per user session, allowing students to resume prior advisory threads. \\ \hline
		\end{longtable}
		
		\subsection{Module FR-2: Dynamic Career Roadmap}
		
		\begin{longtable}{|p{2cm}|p{4cm}|p{9cm}|}
			\hline
			\rowcolor{mainblue!15}\textbf{FR ID} & \textbf{Requirement Name} & \textbf{Detailed Description} \\ \hline
			FR2.1 & Career Role Selection & The system must present a curated catalog of Target Career Roles (e.g., DevOps Engineer, Mobile Developer, Data Engineer, Cloud Architect, Backend Engineer) for student selection. \\ \hline
			FR2.2 & Skill Tree Generation & Upon role selection, the system must automatically generate a hierarchical Skill Tree (Roadmap) displaying technical nodes in prioritized learning sequence, with dependency relationships clearly visualized. \\ \hline
			FR2.3 & Learning Resource Curation & Each skill node must include a minimum of two curated learning resource links sourced from reputable platforms (YouTube, official documentation, Coursera, freeCodeCamp). \\ \hline
			FR2.4 & Progress Tracking & The system must allow students to mark skill nodes as ``Completed'', ``In Progress'', or ``Planned'', and recalculate overall roadmap completion percentage in real time. \\ \hline
			FR2.5 & Adaptive Roadmap Updates & The system must allow roadmap re-generation when the student updates their target career role or when their skill assessment results change significantly ($>$20\% delta). \\ \hline
		\end{longtable}
		
		\subsection{Module FR-3: Skill Gap Analysis}
		
		\begin{longtable}{|p{2cm}|p{4cm}|p{9cm}|}
			\hline
			\rowcolor{mainblue!15}\textbf{FR ID} & \textbf{Requirement Name} & \textbf{Detailed Description} \\ \hline
			FR3.1 & Skill Self-Assessment & The system must provide a structured interface for students to manually input or select their current technical skills from a categorized predefined skill taxonomy. \\ \hline
			FR3.2 & Gap Mapping Analysis & The system must perform a mapping analysis between the student's declared skills and the required skill matrix of the selected Target Career Role, computing gap scores per competency domain. \\ \hline
			FR3.3 & Visual Gap Report & The system must generate a visual report (radar chart and priority list format) identifying missing skills, estimated learning time, and urgency classifications (Critical / Important / Optional). \\ \hline
			FR3.4 & PDF Export & The system must allow students to export the Skill Gap Analysis report as a formatted PDF document. \\ \hline
		\end{longtable}
		
		\subsection{Module FR-4: Market Pulse (Job Intelligence)}
		
		\begin{longtable}{|p{2cm}|p{4cm}|p{9cm}|}
			\hline
			\rowcolor{mainblue!15}\textbf{FR ID} & \textbf{Requirement Name} & \textbf{Detailed Description} \\ \hline
			FR4.1 & Automated Job Data Scraping & The system must automatically scrape data from major IT job portals (LinkedIn, TopCV, ITviec) on a daily scheduled basis via Cron Jobs. \\ \hline
			FR4.2 & NLP Keyword Extraction & The system must perform NLP-based keyword frequency analysis on scraped job descriptions to identify trending technologies, frameworks, and certifications. \\ \hline
			FR4.3 & Interactive Trend Visualization & The system must display interactive trend charts (line charts, bar charts) showing growth or decline in demand for specific IT skills over configurable time windows (30, 90, 180 days). \\ \hline
			FR4.4 & Role-Filtered Market View & The system must allow filtering of market trend data by Target Career Role, enabling students to see demand patterns specific to their chosen specialization. \\ \hline
		\end{longtable}
		
		\subsection{Module FR-5: E-Portfolio Management}
		
		\begin{longtable}{|p{2cm}|p{4cm}|p{9cm}|}
			\hline
			\rowcolor{mainblue!15}\textbf{FR ID} & \textbf{Requirement Name} & \textbf{Detailed Description} \\ \hline
			FR5.1 & GitHub Account Linking & The system must allow students to authenticate and link their GitHub accounts via OAuth 2.0, synchronizing a list of their public repositories. \\ \hline
			FR5.2 & AI README Summarization & The system must use LLM APIs to extract and summarize project objectives, technology stacks, and key features from repository README.md files. \\ \hline
			FR5.3 & Shareable Portfolio URL & The system must generate a unique, publicly accessible URL for each student's E-Portfolio that can be shared with potential employers and interviewers. \\ \hline
			FR5.4 & Portfolio Customization & The system must allow students to reorder displayed projects, add custom descriptions, and toggle project visibility (public/private). \\ \hline
		\end{longtable}
		
		\subsection{Module FR-6: User Account Management}
		
		\begin{longtable}{|p{2cm}|p{4cm}|p{9cm}|}
			\hline
			\rowcolor{mainblue!15}\textbf{FR ID} & \textbf{Requirement Name} & \textbf{Detailed Description} \\ \hline
			FR6.1 & Multi-Provider Authentication & The system must support user authentication via Email/Password combination and Google OAuth 2.0 Single Sign-On. \\ \hline
			FR6.2 & Persistent User State & The system must maintain a persistent database of user chat history, skill assessments, roadmap progress, and portfolio data across login sessions. \\ \hline
			FR6.3 & Role-Based Access Control & The system must enforce RBAC policies distinguishing Student, Career Advisor, Recruiter, and Administrator permission tiers. \\ \hline
			FR6.4 & Profile Management & Students must be able to update their profile information including academic year, institution, target graduation date, and career preferences at any time. \\ \hline
		\end{longtable}
		
		\subsection{Detailed Use Case Specifications}
		
		\subsubsection{UC-01: Student Account Registration}
		
		\begin{longtable}{|p{4.5cm}|p{10.5cm}|}
			\hline
			\rowcolor{gray!20}\textbf{Specification Property} & \textbf{Detailed Specification Content} \\ \hline
			\textbf{Use Case ID} & UC01 \\ \hline
			\textbf{Use Case Name} & Account Registration \\ \hline
			\textbf{Primary Actor} & Student \\ \hline
			\textbf{Pre-conditions} & User has not previously registered. System registration endpoint is online. \\ \hline
			\textbf{Main Flow of Events} &
			1. The student navigates to the registration interface module.\\
			& 2. The system renders the structured registration form.\\
			& 3. The student inputs: Full Name, Institutional Email, Password, University, Academic Year.\\
			& 4. The student triggers the ``Create Account'' action.\\
			& 5. The system validates all field constraints (email format, password strength $\geq$8 chars).\\
			& 6. The system creates a new user database record with hashed credentials.\\
			& 7. The system dispatches a verification token to the provided email address.\\
			& 8. The student activates the account via the email confirmation link.\\ \hline
			\textbf{Alternative Flows} & 3a. Student selects ``Continue with Google''. System initiates OAuth 2.0 flow. Steps 5--7 handled by Google identity provider. \\ \hline
			\textbf{Exception Flows} &
			E1: Input email already exists in the data store --- system rejects transaction with duplicate entry message. \\
			& E2: Password does not meet complexity requirements --- system highlights field with specific error message. \\ \hline
			\textbf{Post-conditions} & A verified user account exists in the system. The student can log in and access onboarding wizard. \\ \hline
		\end{longtable}
		
		\subsubsection{UC-02: User Authentication (Login)}
		
		\begin{longtable}{|p{4.5cm}|p{10.5cm}|}
			\hline
			\rowcolor{gray!20}\textbf{Specification Property} & \textbf{Detailed Specification Content} \\ \hline
			\textbf{Use Case ID} & UC02 \\ \hline
			\textbf{Use Case Name} & User Authentication (Login) \\ \hline
			\textbf{Primary Actors} & Student, Career Advisor, Recruiter, Administrator \\ \hline
			\textbf{Pre-conditions} & User account exists and has been verified. Authentication service is available. \\ \hline
			\textbf{Main Flow of Events} &
			1. The user accesses the centralized authentication interface.\\
			& 2. The user enters identity credentials: Email and Password.\\
			& 3. The user triggers the ``Login'' action.\\
			& 4. The system applies Bcrypt hash comparison against stored credential hash.\\
			& 5. Upon match, the system generates a signed JWT access token (expiry: 24h) and refresh token (expiry: 30d).\\
			& 6. The system loads the user's role-specific dashboard view.\\ \hline
			\textbf{Exception Flows} &
			E1: Credential mismatch --- system displays generic ``Invalid credentials'' error (no specificity to prevent enumeration attacks).\\
			& E2: Account not verified --- system redirects to email confirmation reminder page. \\ \hline
			\textbf{Post-conditions} & User is authenticated and authorized to access role-specific features via valid JWT session. \\ \hline
		\end{longtable}
		
		\subsubsection{UC-03: Generate Personalized Career Roadmap}
		
		\begin{longtable}{|p{4.5cm}|p{10.5cm}|}
			\hline
			\rowcolor{gray!20}\textbf{Specification Property} & \textbf{Detailed Specification Content} \\ \hline
			\textbf{Use Case ID} & UC03 \\ \hline
			\textbf{Use Case Name} & Generate Personalized Career Roadmap \\ \hline
			\textbf{Primary Actor} & Student \\ \hline
			\textbf{Pre-conditions} & Student is authenticated. Student has completed initial skill self-assessment. \\ \hline
			\textbf{Main Flow of Events} &
			1. Student navigates to ``My Roadmap'' module.\\
			& 2. System displays career role selection catalogue.\\
			& 3. Student selects a Target Career Role (e.g., ``DevOps Engineer'').\\
			& 4. System retrieves student's current skill profile from database.\\
			& 5. AI engine computes skill gap delta between student profile and role requirement matrix.\\
			& 6. System renders interactive hierarchical skill tree with prioritized node sequencing.\\
			& 7. Each node displays: skill name, estimated learning hours, prerequisite dependencies, and resource links.\\
			& 8. Student can immediately begin marking nodes and accessing linked resources.\\ \hline
			\textbf{Post-conditions} & Personalized roadmap is stored and accessible from student dashboard at any time. \\ \hline
		\end{longtable}
		
		\section{Non-Functional Requirements (NFR)}
		
		\subsection{Performance Requirements}
		
		\begin{itemize}
			\item \textbf{NFR-P1 Response Latency:} Standard page load and data retrieval operations must complete within 3 seconds under normal network conditions (100 Mbps connection).
			\item \textbf{NFR-P2 AI Response Time:} LLM API-powered responses from the AI Virtual Mentor must return within 10 seconds for standard queries. Complex multi-context queries may extend to 20 seconds with appropriate loading state indication.
			\item \textbf{NFR-P3 Concurrent User Support:} The infrastructure must support a minimum of 500 concurrent active user sessions without measurable degradation of response time SLAs.
			\item \textbf{NFR-P4 Job Scraping Performance:} The automated market data pipeline must complete a full scraping cycle from all configured portals within 2 hours of scheduled trigger time.
		\end{itemize}
		
		\subsection{Security \& Data Protection}
		
		\begin{itemize}
			\item \textbf{NFR-S1 Password Security:} All user passwords must be processed through the Bcrypt hashing algorithm with a minimum cost factor of 12 before persistent storage. Plaintext passwords must never appear in logs or database records.
			\item \textbf{NFR-S2 Transport Security:} All client-server communications must be encrypted via TLS 1.3 protocol minimum. HTTP connections must be automatically redirected to HTTPS.
			\item \textbf{NFR-S3 Access Control:} The system must enforce RBAC at both the API gateway level and individual service layer, preventing privilege escalation attacks.
			\item \textbf{NFR-S4 Input Sanitization:} All user-submitted inputs must be sanitized against SQL injection, XSS (Cross-Site Scripting), and CSRF (Cross-Site Request Forgery) attack vectors.
			\item \textbf{NFR-S5 Data Privacy:} Student academic transcript data and personal profiles must be stored with AES-256 encryption at rest. Data retention policies must comply with Vietnamese personal data protection regulations (Decree 13/2023/ND-CP).
		\end{itemize}
		
		\subsection{Reliability \& Availability}
		
		\begin{itemize}
			\item \textbf{NFR-R1 Service Uptime:} The platform must maintain a minimum Service Level Agreement (SLA) of 99.5\% monthly uptime, permitting no more than 3.6 hours of unscheduled downtime per month.
			\item \textbf{NFR-R2 Fault Tolerance:} The system must implement graceful degradation strategies such that failure of the AI mentoring service does not impact core roadmap viewing or portfolio management functionality.
			\item \textbf{NFR-R3 Data Backup:} Database backups must be automated on a daily schedule with 30-day retention and point-in-time recovery capability within a 1-hour recovery time objective (RTO).
		\end{itemize}
		
		\subsection{Usability Requirements}
		
		\begin{itemize}
			\item \textbf{NFR-U1 Learnability:} New users must be able to complete the initial onboarding workflow (registration, skill assessment, roadmap generation) within 15 minutes without external guidance.
			\item \textbf{NFR-U2 Accessibility:} The web interface must conform to WCAG 2.1 Level AA accessibility standards to accommodate users with visual or motor impairments.
			\item \textbf{NFR-U3 Responsive Design:} The platform must render correctly across all major browsers (Chrome 120+, Firefox 120+, Safari 17+) and device form factors (desktop 1920px, tablet 768px, mobile 375px).
		\end{itemize}
		
		\subsection{Scalability Requirements}
		
		\begin{itemize}
			\item \textbf{NFR-SC1 Horizontal Scaling:} The backend microservices architecture must support horizontal auto-scaling triggered by CPU utilization exceeding 70\% for 3 consecutive minutes.
			\item \textbf{NFR-SC2 Database Scaling:} The primary database must support read replica provisioning to distribute query load as student population scales beyond 5,000 concurrent registered users.
		\end{itemize}
		
		% =====================================
		% CHAPTER III: SYSTEM ARCHITECTURE & UML DESIGN
		% =====================================
		\chapter{System Architecture \& UML Design}
		
		\section{High-Level System Architecture Overview}
		
		The PCOLRP platform adopts a \textbf{Microservices Architecture} pattern deployed on a containerized cloud infrastructure. This architectural decision is motivated by the need to independently scale the AI inference services (which are computationally intensive) from the standard CRUD services (which are I/O-bound). The system decomposes into five primary service domains:
		
		\begin{enumerate}
			\item \textbf{API Gateway Layer:} Centralized ingress point handling authentication token validation, rate limiting, and request routing to downstream microservices.
			\item \textbf{User Identity Service:} Manages authentication, authorization, session management, and RBAC enforcement.
			\item \textbf{Career Intelligence Service:} Orchestrates roadmap generation, skill gap analysis, and AI mentor interactions via LLM API integration.
			\item \textbf{Market Data Service:} Executes scheduled job portal scraping, NLP keyword extraction, and trend analytics computation.
			\item \textbf{Portfolio Service:} Manages GitHub OAuth integration, repository synchronization, AI-powered README summarization, and shareable URL generation.
		\end{enumerate}
		
		\section{System Structural Class Diagram}
		
		\begin{figure}[H]
			\centering
			\resizebox{0.95\textwidth}{!}{
				\begin{tikzpicture}[
					umlclass/.style={
						rectangle split,
						rectangle split parts=3,
						draw=black,
						thick,
						fill=white,
						text width=5cm,
						font=\small,
						align=left,
						inner sep=5pt
					},
					inherits/.style={-{Triangle[open]}, thick},
					assoc/.style={-stealth, thick},
					comp/.style={-{Diamond}, thick}
					]
					
					\node[umlclass] (User) at (9, 6) {
						\textbf{\centerline{User}}
						\nodepart{second}
						userId : Int \\ fullName : String \\ email : String \\ password : String \\ role : String
						\nodepart{third}
						login() \\ logout() \\ updateProfile()
					};
					
					\node[umlclass] (Student) at (1, 0) {
						\textbf{\centerline{Student}}
						\nodepart{second}
						studentId : Int \\ gpaScore : Decimal \\ academicYear : Int \\ institution : String
						\nodepart{third}
						selectCareerRole() \\ submitSkillAssessment() \\ generateRoadmap()
					};
					
					\node[umlclass] (Advisor) at (9, 0) {
						\textbf{\centerline{CareerAdvisor}}
						\nodepart{second}
						advisorId : Int \\ department : String \\ specialization : String
						\nodepart{third}
						viewStudentReport() \\ scheduleSession() \\ addGuidanceNote()
					};
					
					\node[umlclass] (Roadmap) at (1, -7) {
						\textbf{\centerline{Roadmap}}
						\nodepart{second}
						roadmapId : Int \\ targetRole : String \\ completionPct : Float \\ generatedAt : DateTime
						\nodepart{third}
						updateProgress() \\ regenerate() \\ exportPDF()
					};
					
					\node[umlclass] (SkillNode) at (9, -7) {
						\textbf{\centerline{SkillNode}}
						\nodepart{second}
						nodeId : Int \\ skillName : String \\ category : String \\ estimatedHours : Int \\ status : Enum
						\nodepart{third}
						markComplete() \\ getResources() \\ getDependencies()
					};
					
					\node[umlclass] (Portfolio) at (17, 0) {
						\textbf{\centerline{EPortfolio}}
						\nodepart{second}
						portfolioId : Int \\ shareableURL : String \\ githubLogin : String \\ isPublic : Bool
						\nodepart{third}
						syncRepositories() \\ generateSummary() \\ toggleVisibility()
					};
					
					% Inheritance arrows
					\draw[inherits] (Student) -- (User);
					\draw[inherits] (Advisor) -- (User);
					
					% Association arrows
					\draw[assoc] (Student) -- node[left]{\small 1..N} (Roadmap);
					\draw[comp] (Roadmap) -- node[above]{\small 1..N} (SkillNode);
					\draw[assoc] (Student) -- node[above]{\small 1..1} (Portfolio);
					
				\end{tikzpicture}
			}
			\caption{System Core Structural Class Diagram (UML 2.5 Notation)}
		\end{figure}
		
		\section{Use Case Diagram: Full System Actor Interactions}
		
		\begin{figure}[H]
			\centering
			\resizebox{0.92\textwidth}{!}{
				\begin{tikzpicture}[
					actor/.style={draw, circle, minimum size=1cm, fill=blue!10},
					usecase/.style={draw, ellipse, text width=3.5cm, align=center, fill=green!5, minimum height=0.9cm, font=\small},
					system/.style={draw, rectangle, dashed, thick, fill=gray!3},
					arrow/.style={-stealth, thick}
					]
					
					% System boundary
					\node[system, minimum width=14cm, minimum height=12cm, label=above:{\textbf{PCOLRP System Boundary}}] (sys) at (7, -3.5) {};
					
					% Actors
					\node[actor, label=below:{\textbf{Student}}] (student) at (-1, 0) {};
					\node[actor, label=below:{\textbf{Advisor}}] (advisor) at (-1, -5) {};
					\node[actor, label=below:{\textbf{Admin}}] (admin) at (-1, -9) {};
					\node[actor, label=below:{\textbf{Recruiter}}] (recruiter) at (15, -3) {};
					
					% Use Cases
					\node[usecase] (uc1) at (4, 0) {Register \& Login (UC01--02)};
					\node[usecase] (uc2) at (9, 0) {Skill Self-Assessment (UC04)};
					\node[usecase] (uc3) at (4, -3) {Generate Roadmap (UC03)};
					\node[usecase] (uc4) at (9, -3) {AI Virtual Mentor Chat (UC05)};
					\node[usecase] (uc5) at (4, -6) {View Market Pulse (UC06)};
					\node[usecase] (uc6) at (9, -6) {Manage E-Portfolio (UC07)};
					\node[usecase] (uc7) at (6, -9) {View Student Analytics (UC08)};
					
					% Connections
					\draw[arrow] (student) -- (uc1);
					\draw[arrow] (student) -- (uc2);
					\draw[arrow] (student) -- (uc3);
					\draw[arrow] (student) -- (uc4);
					\draw[arrow] (student) -- (uc5);
					\draw[arrow] (student) -- (uc6);
					\draw[arrow] (advisor) -- (uc7);
					\draw[arrow] (advisor) -- (uc4);
					\draw[arrow] (admin) -- (uc1);
					\draw[arrow] (admin) -- (uc7);
					\draw[arrow] (recruiter) -- (uc6);
					
				\end{tikzpicture}
			}
			\caption{System Use Case Diagram --- Actor Interaction Boundary Map}
		\end{figure}
		
		\section{System Sequence Diagrams}
		
		\subsection{Roadmap Generation Sequence Diagram}
		
		\begin{figure}[H]
			\centering
			\resizebox{0.90\textwidth}{!}{
				\begin{tikzpicture}[
					lifeline/.style={draw, thick, dashed},
					actorbox/.style={draw, rectangle, fill=mainblue!15, minimum width=2.5cm, minimum height=0.7cm, font=\small\bfseries, align=center},
					msg/.style={-stealth, thick},
					retmsg/.style={-stealth, dashed, thick}
					]
					
					% Lifeline Headers
					\node[actorbox] (student) at (0, 0) {Student};
					\node[actorbox] (ui) at (3.5, 0) {Frontend UI};
					\node[actorbox] (api) at (7, 0) {API Gateway};
					\node[actorbox] (career) at (10.5, 0) {Career Svc};
					\node[actorbox] (ai) at (14, 0) {AI Engine};
					\node[actorbox] (db) at (17, 0) {Database};
					
					% Lifelines
					\foreach \x in {0, 3.5, 7, 10.5, 14, 17} {
						\draw[lifeline] (\x, -0.35) -- (\x, -10);
					}
					
					% Messages
					\draw[msg] (0,-1) -- node[above, font=\tiny]{selectRole(``DevOps'')} (3.5,-1);
					\draw[msg] (3.5,-1.8) -- node[above, font=\tiny]{POST /roadmap/generate} (7,-1.8);
					\draw[msg] (7,-2.6) -- node[above, font=\tiny]{validateJWT()} (10.5,-2.6);
					\draw[msg] (10.5,-3.4) -- node[above, font=\tiny]{getStudentSkillProfile()} (17,-3.4);
					\draw[retmsg] (17,-4.2) -- node[above, font=\tiny]{skillProfile[]} (10.5,-4.2);
					\draw[msg] (10.5,-5.0) -- node[above, font=\tiny]{computeGapAndGenerate(profile, role)} (14,-5.0);
					\draw[msg] (14,-5.8) -- node[above, font=\tiny]{LLM API: promptRoadmap()} (14,-5.8);
					\draw[retmsg] (14,-6.6) -- node[above, font=\tiny]{roadmapJSON\{nodes, edges\}} (10.5,-6.6);
					\draw[msg] (10.5,-7.4) -- node[above, font=\tiny]{storeRoadmap()} (17,-7.4);
					\draw[retmsg] (17,-8.2) -- node[above, font=\tiny]{roadmapId: 4821} (10.5,-8.2);
					\draw[retmsg] (10.5,-9.0) -- node[above, font=\tiny]{200 OK: roadmapData} (3.5,-9.0);
					\draw[retmsg] (3.5,-9.8) -- node[above, font=\tiny]{renderSkillTree()} (0,-9.8);
					
				\end{tikzpicture}
			}
			\caption{System Sequence Diagram --- Career Roadmap Generation Flow}
		\end{figure}
		
		\section{Data Flow Diagram (DFD)}
		
		\subsection{Context-Level DFD (Level 0)}
		
		\begin{figure}[H]
			\centering
			\resizebox{0.85\textwidth}{!}{
				\begin{tikzpicture}[
					external/.style={rectangle, draw=black, thick, fill=yellow!15, minimum width=2.8cm, minimum height=1.0cm, align=center, font=\small\bfseries},
					process/.style={circle, draw=blue!70, thick, fill=blue!10, minimum size=3.5cm, align=center, font=\small\bfseries},
					datastore/.style={draw, thick, fill=gray!10, minimum width=3cm, minimum height=0.8cm, align=center, font=\small},
					arrow/.style={-stealth, thick}
					]
					
					\node[external] (student_e) at (0, 3) {Student};
					\node[external] (advisor_e) at (0, -3) {Advisor};
					\node[external] (market_e) at (12, 3) {Job Portals};
					\node[external] (recruiter_e) at (12, -3) {Recruiter};
					
					\node[process] (system) at (6, 0) {\textbf{PCOLRP}\\\textbf{Platform}};
					
					\draw[arrow] (student_e) -- node[above, font=\tiny, sloped]{skill data, queries} (system);
					\draw[arrow] (system) -- node[below, font=\tiny, sloped]{roadmap, reports, chat} (student_e);
					\draw[arrow] (advisor_e) -- node[above, font=\tiny, sloped]{guidance notes} (system);
					\draw[arrow] (system) -- node[below, font=\tiny, sloped]{student analytics} (advisor_e);
					\draw[arrow] (market_e) -- node[above, font=\tiny, sloped]{job listing data} (system);
					\draw[arrow] (system) -- node[below, font=\tiny, sloped]{trend insights} (student_e);
					\draw[arrow] (recruiter_e) -- node[above, font=\tiny, sloped]{job postings} (system);
					\draw[arrow] (system) -- node[below, font=\tiny, sloped]{portfolio URL} (recruiter_e);
					
				\end{tikzpicture}
			}
			\caption{System Context-Level Data Flow Diagram (DFD Level 0)}
		\end{figure}
		
		\subsection{Decomposed System DFD (Level 1)}
		
		\begin{figure}[H]
			\centering
			\resizebox{0.95\textwidth}{!}{
				\begin{tikzpicture}[
					process/.style={circle, draw=blue!70, thick, fill=blue!10, minimum size=2.0cm, align=center, font=\tiny\bfseries},
					datastore/.style={draw, thick, fill=gray!10, minimum width=3.5cm, minimum height=0.7cm, align=center, font=\tiny},
					external/.style={rectangle, draw=black, thick, fill=yellow!10, minimum width=2cm, minimum height=0.9cm, align=center, font=\tiny\bfseries},
					arrow/.style={-stealth, thick, font=\tiny}
					]
					
					% External entities
					\node[external] (std) at (0, 8) {Student};
					\node[external] (portal) at (14, 8) {Job Portal};
					
					% Process nodes
					\node[process] (p1) at (2, 5) {P1\\Auth\&Profile};
					\node[process] (p2) at (6, 7) {P2\\Roadmap\\Engine};
					\node[process] (p3) at (10, 7) {P3\\Market\\Pulse};
					\node[process] (p4) at (6, 3) {P4\\AI Mentor};
					\node[process] (p5) at (10, 3) {P5\\Portfolio};
					
					% Data stores
					\node[datastore] (ds1) at (4, 1) {D1: User Profiles DB};
					\node[datastore] (ds2) at (9, 1) {D2: Skill Nodes DB};
					\node[datastore] (ds3) at (14, 4) {D3: Market Jobs DB};
					
					% Flow connections
					\draw[arrow] (std) -- (p1);
					\draw[arrow] (p1) -- (p2);
					\draw[arrow] (p1) -- (p4);
					\draw[arrow] (p1) -- (p5);
					\draw[arrow] (portal) -- (p3);
					\draw[arrow] (p3) -- (ds3);
					\draw[arrow] (ds3) -- (p2);
					\draw[arrow] (p2) -- (ds2);
					\draw[arrow] (p1) -- (ds1);
					\draw[arrow] (ds1) -- (p4);
					\draw[arrow] (ds2) -- (p4);
					
				\end{tikzpicture}
			}
			\caption{Decomposed Operational Data Routing Blueprint (DFD Level 1)}
		\end{figure}
		
		\subsection{Process Optimization Swimlane Mapping for AI Engines}
		
		\begin{figure}[H]
			\centering
			\resizebox{0.95\textwidth}{!}{
				\begin{tikzpicture}[
					swimlane/.style={draw=black, thick, minimum width=4.5cm, minimum height=6cm, fill=gray!5, align=center},
					node distance=0.5cm
					]
					\node[swimlane] (lane1) at (0,0) {};
					\node[swimlane] (lane2) at (4.6,0) {};
					\node[swimlane] (lane3) at (9.2,0) {};
					
					\node[above=0.1cm of lane1.north, font=\bfseries] {Student Client UI};
					\node[above=0.1cm of lane2.north, font=\bfseries] {Backend Service Core};
					\node[above=0.1cm of lane3.north, font=\bfseries] {AI Engine Process};
					
					\node[draw, rectangle, rounded corners, fill=orange!15, text width=3.5cm, align=center] (step1) at (0, 1.5) {Target Career Node Selection};
					\node[draw, rectangle, fill=blue!15, text width=3.5cm, align=center] (step2) at (4.6, 0) {Aggregate Profile Context Data};
					\node[draw, diamond, aspect=2, fill=green!15, text width=2.5cm, align=center, inner sep=1pt] (step3) at (9.2, -1) {Compute Technology Gap};
					\node[draw, rectangle, fill=purple!15, text width=3.5cm, align=center] (step4) at (9.2, -3) {Generate LLM Roadmap Prompt};
					\node[draw, rectangle, rounded corners, fill=yellow!15, text width=3.5cm, align=center] (step5) at (0, -2.5) {Render Skill Tree Visualization};
					
					\draw[-stealth, thick] (step1) -- (step2);
					\draw[-stealth, thick] (step2) -- (step3);
					\draw[-stealth, thick] (step3) -- node[right, font=\tiny]{Gap $>$ threshold} (step4);
					\draw[-stealth, thick] (step4) -| (step5);
				\end{tikzpicture}
			}
			\caption{Cross-Functional Process Optimization Flowchart for Artificial Intelligence Pipeline}
		\end{figure}
		
		% =====================================
		% CHAPTER IV: DATABASE & INTERFACE DESIGN
		% =====================================
		\chapter{Database \& Interface Design}
		
		\section{Conceptual Data Modeling (Conceptual ERD)}
		
		\begin{figure}[htbp]
			\begin{center}
				\begin{singlespace}
					\begin{tikzpicture}[
						entity/.style={rectangle, draw=blue!70, fill=blue!5, thick, minimum width=3.2cm, minimum height=0.9cm, rounded corners=2pt},
						relationship/.style={diamond, draw=red!70, fill=red!5, thick, minimum width=2.2cm, minimum height=0.9cm, aspect=1.5},
						every node/.style={font=\small\bfseries},
						node distance=1.4cm and 1.2cm
						]
						
						% Tier 1: Core orientation management
						\node[entity] (student) {Student Profile};
						\node[relationship] (select) [right=of student] {Selects};
						\node[entity] (path) [right=of select] {Tech Path};
						
						% Tier 2: Middle connection
						\node[relationship] (have) [below=of student] {Participates};
						\node[relationship] (contain) [below=of path] {Contains};
						
						% Tier 3: Detail nodes
						\node[entity] (mentor) [below=of have] {Mentor Session};
						\node[entity] (skill) [below=of contain] {Skill Node};
						
						% Tier 4: Branch from Skill Node
						\node[relationship] (link) [below left=of skill] {Links};
						\node[entity] (course) [below=of link] {Course Repo};
						
						\node[relationship] (update) [below right=of skill] {Updates};
						\node[entity] (job) [below=of update] {Job Trend};
						
						\node[relationship] (owns) [right=of student, xshift=7cm] {Owns};
						\node[entity] (portfolio) [right=of owns] {E-Portfolio};
						
						% Cardinality lines
						\path[draw, thick] (student) -- node[above, near start] {N} (select);
						\path[draw, thick] (select) -- node[above, near end] {M} (path);
						
						\path[draw, thick] (student) -- node[left, near start] {1} (have);
						\path[draw, thick] (have) -- node[left, near end] {N} (mentor);
						
						\path[draw, thick] (path) -- node[left, near start] {N} (contain);
						\path[draw, thick] (contain) -- node[left, near end] {M} (skill);
						
						\path[draw, thick] (skill) -- node[above left, near start] {1} (link);
						\path[draw, thick] (link) -- node[below left, near end] {N} (course);
						
						\path[draw, thick] (skill) -- node[above right, near start] {N} (update);
						\path[draw, thick] (update) -- node[below right, near end] {M} (job);
						
						\path[draw, thick] (student) -- node[above] {1} (owns);
						\path[draw, thick] (owns) -- node[above] {1} (portfolio);
						
					\end{tikzpicture}
				\end{singlespace}
			\end{center}
			\caption{System Core Data Framework Architecture (Conceptual ERD)}
		\end{figure}
		
		\section{Logical Data Schema Organization (Logical ERD)}
		
		The following table documents all inter-table relationships in the logical data model, mapped from the conceptual ERD entities to normalized relational database tables:
		
		\begin{longtable}{|p{3.5cm}|p{3.0cm}|p{3.5cm}|p{4.5cm}|}
			\hline
			\rowcolor{gray!20}\textbf{Source Table} & \textbf{Relationship} & \textbf{Target Table} & \textbf{Join Attribute} \\ \hline
			StudentProfiles & 1 -- N & Enrollments & StudentID FK \\ \hline
			Courses & 1 -- N & Enrollments & CourseID FK \\ \hline
			StudentProfiles & 1 -- N & StudentSkills & StudentID FK \\ \hline
			Skills & 1 -- N & StudentSkills & SkillID FK \\ \hline
			Courses & N -- M & CourseSkills & CourseID, SkillID FK \\ \hline
			Skills & 1 -- N & CourseSkills & SkillID FK \\ \hline
			StudentProfiles & 1 -- N & RoadmapProgress & StudentID FK \\ \hline
			SkillNodes & 1 -- N & RoadmapProgress & NodeID FK \\ \hline
			StudentProfiles & 1 -- 1 & EPortfolios & StudentID FK \\ \hline
			EPortfolios & 1 -- N & Repositories & PortfolioID FK \\ \hline
			JobPostings & N -- M & Skills & via JobSkillsTags \\ \hline
		\end{longtable}
		
		\section{Detailed Data Dictionary}
		
		\subsection{Table: USERS (Student Core Identity Store)}
		
		\begin{tabular}{|p{3.5cm}|p{3.2cm}|p{2.5cm}|p{5.8cm}|}
			\hline
			\rowcolor{mainblue!10}\textbf{Field Attribute} & \textbf{Data Type} & \textbf{Key / Constraint} & \textbf{Functional Description} \\ \hline
			id & INT & PK, Auto Inc & Unique structural identifier for each system user. \\ \hline
			username & VARCHAR(50) & Unique, Not Null & Unique string used for secure user login. \\ \hline
			email & VARCHAR(100) & Unique, Not Null & Primary institutional email address. \\ \hline
			password\_hash & VARCHAR(255) & Not Null & Salted Bcrypt cryptographic hash of user credential. \\ \hline
			role & VARCHAR(20) & Not Null & Privileges tier: `Student', `Admin', `Mentor', `Recruiter'. \\ \hline
			gpa\_score & DECIMAL(3,2) & Nullable & Academic performance grade point average (0.00--4.00). \\ \hline
			github\_url & VARCHAR(255) & Nullable & Linked public GitHub profile URL. \\ \hline
			created\_at & DATETIME & Not Null & Account creation timestamp (UTC). \\ \hline
			is\_verified & BOOLEAN & Default FALSE & Email verification confirmation flag. \\ \hline
		\end{tabular}
		
		\subsection{Table: SKILL\_NODES (Technology Competency Taxonomy)}
		
		\begin{tabular}{|p{3.5cm}|p{3.2cm}|p{2.5cm}|p{5.8cm}|}
			\hline
			\rowcolor{mainblue!10}\textbf{Field Attribute} & \textbf{Data Type} & \textbf{Key / Constraint} & \textbf{Functional Description} \\ \hline
			node\_id & INT & PK, Auto Inc & Unique identifier for each skill node. \\ \hline
			skill\_name & VARCHAR(100) & Not Null & Canonical name of the technical skill (e.g., ``Docker''). \\ \hline
			category & VARCHAR(50) & Not Null & Domain classification: DevOps, AI/ML, Frontend, etc. \\ \hline
			career\_role & VARCHAR(50) & Not Null & Target role this node belongs to. \\ \hline
			priority\_rank & INT & Not Null & Learning sequence priority (lower = learn first). \\ \hline
			estimated\_hours & INT & Not Null & Estimated learning hours to reach proficiency. \\ \hline
			prerequisite\_ids & JSON & Nullable & Array of node\_ids that must be completed first. \\ \hline
		\end{tabular}
		
		\subsection{Table: ROADMAP\_PROGRESS (Individual Learning Tracking)}
		
		\begin{tabular}{|p{3.5cm}|p{3.2cm}|p{2.5cm}|p{5.8cm}|}
			\hline
			\rowcolor{mainblue!10}\textbf{Field Attribute} & \textbf{Data Type} & \textbf{Key / Constraint} & \textbf{Functional Description} \\ \hline
			progress\_id & INT & PK, Auto Inc & Unique record identifier. \\ \hline
			student\_id & INT & FK $\to$ USERS & References the owning student account. \\ \hline
			node\_id & INT & FK $\to$ SKILL\_NODES & References the associated skill node. \\ \hline
			status & ENUM & Not Null & `Planned', `In Progress', `Completed'. \\ \hline
			completed\_at & DATETIME & Nullable & Timestamp when node was marked completed. \\ \hline
			notes & TEXT & Nullable & Optional student notes for the learning node. \\ \hline
		\end{tabular}
		
		\subsection{Table: MARKET\_PULSE\_JOBS (Job Market Intelligence Cache)}
		
		\begin{tabular}{|p{3.5cm}|p{3.2cm}|p{2.5cm}|p{5.8cm}|}
			\hline
			\rowcolor{mainblue!10}\textbf{Field Attribute} & \textbf{Data Type} & \textbf{Key / Constraint} & \textbf{Functional Description} \\ \hline
			job\_id & INT & PK, Auto Inc & Unique job listing identifier. \\ \hline
			job\_title & VARCHAR(200) & Not Null & Extracted job position title. \\ \hline
			company\_name & VARCHAR(150) & Not Null & Employer company name. \\ \hline
			source\_portal & VARCHAR(50) & Not Null & Origin portal: `LinkedIn', `TopCV', `ITviec'. \\ \hline
			scraped\_at & DATETIME & Not Null & UTC timestamp of data collection. \\ \hline
			extracted\_skills & JSON & Not Null & NLP-extracted technology keywords array. \\ \hline
			career\_category & VARCHAR(50) & Not Null & Classified role domain for filtering. \\ \hline
		\end{tabular}
		
		\section{System User Interface Design Wireframes}
		
		\subsection{Dashboard Overview Screen}
		
		The main student dashboard provides a unified command center consolidating four information widgets:
		
		\begin{itemize}
			\item \textbf{Roadmap Progress Widget:} Circular progress indicator showing overall completion percentage of the active learning roadmap, with the next recommended skill node highlighted.
			\item \textbf{Market Pulse Snapshot:} A mini bar chart showing the top 5 trending skills in the student's target career domain over the past 30 days.
			\item \textbf{AI Mentor Quick Access:} A persistent floating action button enabling immediate access to the conversational AI mentor from any screen.
			\item \textbf{Portfolio Strength Score:} A composite score (0--100) evaluating portfolio coherence, GitHub activity, and roadmap completion alignment.
		\end{itemize}
		
		\subsection{Skill Tree Visualization Screen}
		
		The roadmap visualization screen renders an interactive directed acyclic graph (DAG) where:
		
		\begin{itemize}
			\item Nodes are color-coded by status: \textbf{\color{accentgreen}Green} (Completed), \textbf{\color{accentorange}Orange} (In Progress), \textbf{Gray} (Planned), \textbf{\color{red}Red} (Critical Gap).
			\item Edge lines represent prerequisite dependencies between skill nodes.
			\item Clicking any node expands a side panel showing skill description, learning resources, estimated hours, and market demand data.
			\item Zoom and pan gestures allow navigation of large skill trees (DevOps path averages 42 nodes).
		\end{itemize}
		
		% =====================================
		% CHAPTER V: IMPLEMENTATION & TESTING
		% =====================================
		\chapter{Implementation, AI Integration \& Testing}
		
		\section{Technology Stack Selection}
		
		\subsection{Frontend Technology}
		
		\begin{longtable}{|p{3cm}|p{4cm}|p{8cm}|}
			\hline
			\rowcolor{gray!20}\textbf{Component} & \textbf{Technology} & \textbf{Justification} \\ \hline
			Framework & React 18 + TypeScript & Component-driven architecture with strong type safety for complex interactive UI (skill trees, charts). \\ \hline
			Routing & React Router v6 & Declarative client-side routing with nested layout support. \\ \hline
			State Management & Redux Toolkit & Centralized state management for cross-component roadmap progress synchronization. \\ \hline
			Graph Visualization & React Flow / D3.js & Interactive node-edge rendering for skill tree DAG visualization. \\ \hline
			UI Component Library & Tailwind CSS + shadcn/ui & Utility-first styling enabling rapid, consistent component development. \\ \hline
			Chart Library & Recharts & Declarative React-native charting for Market Pulse trend visualizations. \\ \hline
		\end{longtable}
		
		\subsection{Backend Technology}
		
		\begin{longtable}{|p{3cm}|p{4cm}|p{8cm}|}
			\hline
			\rowcolor{gray!20}\textbf{Component} & \textbf{Technology} & \textbf{Justification} \\ \hline
			Runtime & Node.js 20 LTS & Non-blocking I/O model well-suited for concurrent API request handling and LLM streaming. \\ \hline
			Framework & Express.js 4.x & Minimal, flexible HTTP server framework with extensive middleware ecosystem. \\ \hline
			Authentication & Passport.js + JWT & Modular authentication supporting multiple OAuth strategies. \\ \hline
			ORM & Prisma 5.x & Type-safe database client with migration management and schema introspection. \\ \hline
			Task Scheduling & node-cron & Lightweight cron job executor for automated market data scraping pipelines. \\ \hline
			AI Integration & OpenAI SDK v4 & Official Node.js client for GPT-4o API with streaming support. \\ \hline
		\end{longtable}
		
		\subsection{Infrastructure \& DevOps}
		
		\begin{longtable}{|p{3cm}|p{4cm}|p{8cm}|}
			\hline
			\rowcolor{gray!20}\textbf{Component} & \textbf{Technology} & \textbf{Justification} \\ \hline
			Database & PostgreSQL 16 & ACID-compliant relational database with strong JSON support for dynamic skill taxonomies. \\ \hline
			Caching & Redis 7 & In-memory caching layer for LLM response memoization and session management. \\ \hline
			Containerization & Docker + Docker Compose & Reproducible environment packaging for development and production parity. \\ \hline
			CI/CD & GitHub Actions & Automated testing, linting, and deployment pipeline on each pull request merge. \\ \hline
			Hosting & Railway / Render (pilot) & Cloud PaaS simplifying Node.js deployment for the UTH pilot phase. \\ \hline
		\end{longtable}
		
		\section{Physical Relational Database Implementation (SQL DDL)}
		
		\begin{lstlisting}[language=SQL, caption={Core Database Schema DDL --- PCOLRP PostgreSQL Implementation}]
			-- 1. Create High-Quality Student Profile Schema
			CREATE TABLE StudentProfiles (
			StudentID    SERIAL PRIMARY KEY,
			FullName     VARCHAR(100) NOT NULL,
			Email        VARCHAR(100) UNIQUE NOT NULL,
			PasswordHash VARCHAR(255) NOT NULL,
			AcademicYear INT          CHECK (AcademicYear BETWEEN 1 AND 5),
			Specialization VARCHAR(50) DEFAULT 'Software Engineering',
			GPAScore     DECIMAL(3,2) CHECK (GPAScore BETWEEN 0.0 AND 4.0),
			GitHubURL    VARCHAR(255) NULL,
			IsVerified   BOOLEAN      DEFAULT FALSE,
			Role         VARCHAR(20)  DEFAULT 'Student',
			CreatedAt    TIMESTAMPTZ  DEFAULT NOW()
			);
			
			-- 2. Create Technology Skill Nodes Taxonomy Store
			CREATE TABLE SkillNodes (
			NodeID          SERIAL PRIMARY KEY,
			SkillName       VARCHAR(100) UNIQUE NOT NULL,
			Category        VARCHAR(50)  NOT NULL,
			CareerRole      VARCHAR(50)  NOT NULL,
			PriorityRank    INT          NOT NULL,
			EstimatedHours  INT          NOT NULL DEFAULT 10,
			PrerequisiteIDs JSONB        DEFAULT '[]'
			);
			
			-- 3. Create Student Learning Progress Tracking Matrix
			CREATE TABLE RoadmapProgress (
			ProgressID  SERIAL PRIMARY KEY,
			StudentID   INT REFERENCES StudentProfiles(StudentID) ON DELETE CASCADE,
			NodeID      INT REFERENCES SkillNodes(NodeID),
			Status      VARCHAR(20) DEFAULT 'Planned'
			CHECK (Status IN ('Planned', 'In Progress', 'Completed')),
			CompletedAt TIMESTAMPTZ NULL,
			Notes       TEXT,
			UpdatedAt   TIMESTAMPTZ DEFAULT NOW(),
			UNIQUE (StudentID, NodeID)
			);
			
			-- 4. Create Market Pulse Job Intelligence Cache
			CREATE TABLE MarketPulseJobs (
			JobID           SERIAL PRIMARY KEY,
			JobTitle        VARCHAR(200) NOT NULL,
			CompanyName     VARCHAR(150) NOT NULL,
			SourcePortal    VARCHAR(50)  NOT NULL,
			ExtractedSkills JSONB        DEFAULT '[]',
			CareerCategory  VARCHAR(50),
			ScrapedAt       TIMESTAMPTZ  DEFAULT NOW()
			);
			
			-- 5. Create E-Portfolio Entity
			CREATE TABLE EPortfolios (
			PortfolioID  SERIAL PRIMARY KEY,
			StudentID    INT UNIQUE REFERENCES StudentProfiles(StudentID) ON DELETE CASCADE,
			ShareableURL VARCHAR(255) UNIQUE NOT NULL,
			GitHubLogin  VARCHAR(100),
			IsPublic     BOOLEAN DEFAULT TRUE,
			LastSynced   TIMESTAMPTZ
			);
			
			-- 6. Create Repository Listings Cache
			CREATE TABLE Repositories (
			RepoID        SERIAL PRIMARY KEY,
			PortfolioID   INT REFERENCES EPortfolios(PortfolioID) ON DELETE CASCADE,
			RepoName      VARCHAR(200) NOT NULL,
			RepoURL       VARCHAR(255) NOT NULL,
			Description   TEXT,
			AISummary     TEXT,
			TechStack     JSONB DEFAULT '[]',
			StarCount     INT DEFAULT 0,
			IsVisible     BOOLEAN DEFAULT TRUE
			);
		\end{lstlisting}
		
		\section{AI Integration Architecture}
		
		\subsection{LLM API Integration Specification}
		
		The AI Virtual Mentor module integrates with the OpenAI GPT-4o API using a structured prompt engineering approach. Each API call constructs a system prompt that injects the student's skill profile, academic transcript summary, and target career role context to ensure highly personalized, grounded responses.
		
		\begin{lstlisting}[
			caption={AI Mentor Prompt Engineering Module (Node.js)},
			basicstyle=\ttfamily\tiny]
			async function generateMentorResponse(studentContext, userQuery) {
				const systemPrompt = `
				You are an expert Software Engineering career advisor
				for a Vietnamese university student with the following profile:
				- Target Career Role: ${studentContext.targetRole}
				- Current GPA: ${studentContext.gpa}
				- Completed Skills: ${studentContext.completedSkills.join(', ')}
				- Critical Skill Gaps: ${studentContext.skillGaps.join(', ')}
				- Academic Year: Year ${studentContext.academicYear}
				
				Provide specific, actionable career advice tailored to this profile.
				Reference specific technologies, project ideas, and resources relevant
				to the Vietnamese IT job market where appropriate.
				Respond in ${studentContext.preferredLanguage}.
				`;
				
				const response = await openai.chat.completions.create({
					model: "gpt-4o",
					messages: [
					{ role: "system", content: systemPrompt },
					...studentContext.conversationHistory,
					{ role: "user", content: userQuery }
					],
					max_tokens: 1000,
					temperature: 0.7,
					stream: true
				});
				
				return response; // Streamed to client via SSE
			}
		\end{lstlisting}
		
		\subsection{Market Data Scraping Pipeline}
		
		\begin{figure}[H]
			\centering
			\resizebox{0.80\textwidth}{!}{
				\begin{tikzpicture}[
					block/.style={rectangle, draw=blue!80, fill=blue!10, text width=7cm, text centered, rounded corners, minimum height=0.9cm, font=\small},
					dec/.style={diamond, draw=orange!80, fill=orange!10, text width=2.5cm, text centered, aspect=2, font=\small},
					arrow/.style={thick, ->, >=stealth, draw=gray!80}
					]
					
					\node[block] (s) {Cron Job triggers automatically at 00:00 UTC daily};
					\node[block] (p1) [below=0.4cm of s] {Dispatches HTTP requests to LinkedIn, TopCV, ITviec API endpoints};
					\node[block] (p2) [below=0.4cm of p1] {Parses HTML/JSON payloads: Job Titles, Descriptions, Company Names};
					\node[dec] (d1) [below=0.6cm of p2] {Duplicate record?};
					\node[block] (p3) [below=0.8cm of d1] {Routes text payload into spaCy NLP Entity Extraction pipeline};
					\node[block] (p4) [below=0.4cm of p3] {Extracts skill keywords and classifies career category};
					\node[block] (p5) [below=0.4cm of p4] {Persists normalized records to MarketPulseJobs table (PostgreSQL)};
					\node[block] (p6) [below=0.4cm of p5] {Invalidates Redis trend cache --- forces chart re-computation on next request};
					
					\draw[arrow] (s) -- (p1);
					\draw[arrow] (p1) -- (p2);
					\draw[arrow] (p2) -- (d1);
					\draw[arrow] (d1) -- node[right, font=\tiny] {No --- new record} (p3);
					\draw[arrow] (d1.east) -- ++(2.5,0) node[right, font=\tiny] {Yes --- skip} |- (p5);
					\draw[arrow] (p3) -- (p4);
					\draw[arrow] (p4) -- (p5);
					\draw[arrow] (p5) -- (p6);
					
				\end{tikzpicture}
			}
			\caption{Automated Market Data Synchronization Workflow (Activity Diagram)}
		\end{figure}
		
		\section{System Testing Strategy}
		
		\subsection{Testing Framework Overview}
		
		The validation lifecycle is organized across four integration layers, each targeting a distinct scope of the system's operational surface:
		
		\begin{longtable}{|p{3cm}|p{3cm}|p{4cm}|p{5cm}|}
			\hline
			\rowcolor{gray!20}\textbf{Test Layer} & \textbf{Tool / Framework} & \textbf{Coverage Target} & \textbf{Success Criteria} \\ \hline
			Unit Testing & Jest + Supertest & Service functions, utility modules, data validators & $\geq$85\% line coverage \\ \hline
			Integration Testing & Jest + Testcontainers & API endpoint to database flows, AI service mocking & All endpoints return correct HTTP status codes \\ \hline
			System / Load Testing & k6 (Grafana) & 500 concurrent virtual users, 10-minute sustained load & p95 response time $<$ 3s, error rate $<$ 1\% \\ \hline
			UAT & Manual + UTH students & Real-world onboarding workflow, roadmap usability & $\geq$80\% task completion rate on 20-user cohort \\ \hline
		\end{longtable}
		
		\subsection{Sample Unit Test Cases}
		
		\begin{longtable}{|p{1.5cm}|p{4cm}|p{4.5cm}|p{5cm}|}
			\hline
			\rowcolor{gray!20}\textbf{TC ID} & \textbf{Test Case Name} & \textbf{Test Input} & \textbf{Expected Output} \\ \hline
			TC-U01 & Valid user registration & email: test@uth.edu.vn, password: SecurePass123! & HTTP 201, userId returned in body \\ \hline
			TC-U02 & Duplicate email registration & Existing email address repeated & HTTP 409 Conflict with error message \\ \hline
			TC-U03 & Login with correct credentials & Valid email + password pair & HTTP 200, JWT token in response body \\ \hline
			TC-U04 & Login with wrong password & Valid email + incorrect password & HTTP 401 Unauthorized \\ \hline
			TC-U05 & Roadmap generation for DevOps & studentId: 1, role: ``DevOps Engineer'' & HTTP 200, roadmapId created, nodes $\geq$ 30 \\ \hline
			TC-U06 & Skill node status update & nodeId: 15, status: ``Completed'' & HTTP 200, completedAt timestamp set \\ \hline
			TC-U07 & AI Mentor response generation & Valid student context + query string & HTTP 200, response.content length $>$ 50 chars \\ \hline
			TC-U08 & Portfolio URL generation & studentId: 3, portfolioData valid & HTTP 201, unique shareableURL returned \\ \hline
		\end{longtable}
		
		\subsection{System Performance Test Results (Simulated)}
		
		The following summarizes the key performance metrics obtained during load testing simulation using k6 with a 500 virtual user ramp-up profile over a 10-minute duration:
		
		\begin{longtable}{|p{5cm}|p{3cm}|p{3cm}|p{4cm}|}
			\hline
			\rowcolor{gray!20}\textbf{Metric} & \textbf{Target SLA} & \textbf{Measured Value} & \textbf{Status} \\ \hline
			Average Response Time (API) & $<$ 3,000 ms & 842 ms & \textbf{\color{accentgreen}PASS} \\ \hline
			95th Percentile Response Time & $<$ 3,000 ms & 1,847 ms & \textbf{\color{accentgreen}PASS} \\ \hline
			AI Mentor P95 Response & $<$ 10,000 ms & 7,234 ms & \textbf{\color{accentgreen}PASS} \\ \hline
			HTTP Error Rate & $<$ 1\% & 0.23\% & \textbf{\color{accentgreen}PASS} \\ \hline
			Concurrent User Throughput & 500 users & 512 users sustained & \textbf{\color{accentgreen}PASS} \\ \hline
			Database Query P95 & $<$ 500 ms & 187 ms & \textbf{\color{accentgreen}PASS} \\ \hline
		\end{longtable}
		
		% =====================================
		% CHAPTER VI: CONCLUSION & FUTURE WORK
		% =====================================
		\chapter{Conclusion \& Future Work}
		
		\section{Summary of Achieved Deliverables}
		
		This project has successfully designed, specified, and implemented the foundational architecture for the \textbf{Personalized Career Orientation \& Learning Roadmap Platform (PCOLRP)} --- an intelligent, AI-augmented web platform purpose-built to bridge the widening gap between academic software engineering education and modern industry career expectations for students at the University of Transport Ho Chi Minh City.
		
		The six core deliverables achieved across the project lifecycle are:
		
		\begin{enumerate}
			\item \textbf{Comprehensive SRS Documentation:} A complete Software Requirements Specification in IEEE 830-1998 format, covering 24 functional requirements across 6 modular subsystems and 5 non-functional requirement domains (Performance, Security, Reliability, Usability, Scalability).
			
			\item \textbf{Full UML Design Suite:} A complete UML design artifact package comprising: Structural Class Diagram (7 entities), System Use Case Diagram (8 use cases, 4 actors), Sequence Diagrams for primary workflows, Context-Level and Level-1 Data Flow Diagrams, and a Cross-Functional AI Pipeline Swimlane Map.
			
			\item \textbf{Relational Database Architecture:} A fully normalized PostgreSQL relational schema implementing 8 core tables with referential integrity constraints, JSONB dynamic attribute support, and indexing strategies optimized for roadmap query performance.
			
			\item \textbf{AI Integration Architecture:} A production-ready LLM integration module using GPT-4o with prompt engineering techniques that inject student profile context (GPA, skills, career role, conversation history) to deliver genuinely personalized career mentoring at scale.
			
			\item \textbf{Automated Market Intelligence Pipeline:} A scheduled data scraping and NLP extraction pipeline that continuously harvests job market demand signals from three major Vietnamese IT job portals, enabling real-time alignment of learning recommendations with employer expectations.
			
			\item \textbf{Validated Testing Strategy:} A multi-layer testing framework spanning Unit, Integration, System Load, and User Acceptance Testing, with documented test cases and simulated performance results confirming all SLA targets are achievable under the projected UTH pilot load profile.
		\end{enumerate}
		
		\section{Academic and Practical Impact Assessment}
		
		\subsection{Impact on Students}
		
		The platform directly addresses the three core problem clusters identified in Chapter~I:
		
		\begin{itemize}
			\item \textbf{Curriculum Gap Resolution:} By continuously aggregating real-time job market keyword data, the platform's skill recommendations automatically reflect current employer demands rather than lagging academic curricula, reducing the skill gap exposure for graduating students.
			\item \textbf{Choice Paralysis Elimination:} The AI-driven personalized roadmap generation replaces the overwhelming landscape of generic online learning paths with a single, tailored, priority-sequenced learning path grounded in the student's actual competency baseline.
			\item \textbf{Portfolio Coherence:} The E-Portfolio system connects GitHub repositories to the student's skill roadmap narrative, enabling students to present a coherent, employer-optimized professional story rather than a disconnected collection of class projects.
		\end{itemize}
		
		\subsection{Impact on Career Advisors}
		
		Career advisors gain access to a data-rich analytics dashboard that surfaces aggregate skill gap patterns across student cohorts, enabling evidence-based intervention strategies and more effective advisory session preparation. The AI mentoring layer handles high-volume routine career queries, freeing advisor capacity for deeper, high-value consultation engagements.
		
		\section{Constraints and Limitations}
		
		\begin{enumerate}
			\item \textbf{LLM API Cost Dependency:} The AI Virtual Mentor subsystem incurs variable operational costs proportional to token consumption per student interaction. At scale, LLM API billing requires budget planning and potential cost-optimization strategies (response caching, smaller model fallback routing).
			
			\item \textbf{Scraping Compliance Risk:} Automated job portal scraping operates in a legally ambiguous zone with respect to Terms of Service compliance across LinkedIn and other commercial platforms. Production deployment requires legal review and potential pivot to official API partnerships or data licensing agreements.
			
			\item \textbf{Skill Taxonomy Maintenance:} The static skill taxonomy curated for the initial role catalog requires ongoing editorial maintenance as new technologies emerge and others deprecate. An automated community-driven taxonomy update mechanism is required for long-term operational sustainability.
			
			\item \textbf{Cold-Start Problem:} New students without established GitHub activity or completed skill assessments receive less personalized AI recommendations until sufficient behavioral data accumulates. Onboarding wizard design must compensate with richer initial questionnaire data collection.
		\end{enumerate}
		
		\section{Future Research Directions and Roadmap}
		
		The following enhancement vectors are planned for subsequent development phases, organized by strategic priority:
		
		\subsection{Phase 2: Deep Learning Recommendation Engine (Sprint 5--8)}
		
		\begin{itemize}
			\item Replace the current LLM-API-dependent roadmap generation with a locally trained Collaborative Filtering recommendation model trained on anonymized aggregate student progress data. This approach reduces operational costs and improves recommendation latency to sub-second response times.
			\item Implement a Recurrent Neural Network (LSTM-based) time-series model for job market demand forecasting, enabling the platform to surface ``emerging'' skills before they reach peak market saturation --- giving students a competitive advantage.
		\end{itemize}
		
		\subsection{Phase 3: Multi-Institution Expansion (Sprint 9--12)}
		
		\begin{itemize}
			\item Extend the platform's multi-tenancy architecture to support additional Vietnamese universities (HCMUT, FPT University, UIT) with institution-specific curriculum integration and cohort-level analytics dashboards for department heads.
			\item Develop a REST API layer enabling third-party academic management systems (LMS platforms such as Moodle) to synchronize course completion data directly into student skill profiles.
		\end{itemize}
		
		\subsection{Phase 4: Mobile-Native Applications (Sprint 13--16)}
		
		\begin{itemize}
			\item Develop native iOS and Android companion applications using React Native (code-sharing with the existing React web frontend) to enable on-the-go skill node progress tracking, AI mentor chat, and push notification reminders for learning milestones.
			\item Implement offline-capable skill tree viewing with background sync for students with intermittent connectivity in remote areas.
		\end{itemize}
		
		\subsection{Phase 5: Employer Integration Ecosystem (Sprint 17--20)}
		
		\begin{itemize}
			\item Build a recruiter-facing portal enabling companies to publish targeted skill-matching job postings that the system automatically maps to student profiles meeting the specified competency requirements.
			\item Develop an anonymized talent pool API allowing partner companies to request ranked candidate shortlists based on skill compatibility scores, accelerating the internship and graduate recruitment pipeline for UTH students.
		\end{itemize}
		
		\section{Reflections on the Software Engineering Process}
		
		The development of this project under Pressman's Software Engineering Framework (7th Edition) provided the team with tangible experience in the disciplined application of software engineering practices often underemphasized in project-based learning environments:
		
		\begin{itemize}
			\item \textbf{Requirements Traceability:} The structured SRS process revealed three functional requirements during stakeholder review that were initially omitted from the use case workshops (FR3.4 PDF Export, FR4.4 Role-Filtered Market View, FR5.4 Portfolio Customization), demonstrating the value of formal requirements elicitation over ad-hoc feature ideation.
			\item \textbf{Architectural Decision Documentation:} Formally documenting the rationale for microservices vs. monolith architecture selection early in the project prevented scope creep and enabled consistent technical decision-making across all team members throughout the implementation phase.
			\item \textbf{Testing as First-Class Activity:} Treating the testing strategy as a parallel planning artifact (not an afterthought) enabled the team to design testable interfaces from the first implementation sprint, significantly reducing defect discovery latency.
		\end{itemize}
		
		\vspace{1cm}
		
		\begin{tcolorbox}[colback=mainblue!5, colframe=mainblue, title={\textbf{Final Conclusion}}]
			The PCOLRP platform represents a meaningful contribution to the Vietnamese Software Engineering education ecosystem. By intelligently converging personalized AI mentoring, real-time market intelligence, and structured skill development guidance into a unified student-centric platform, this project demonstrates that the gap between academic preparation and professional readiness can be systematically narrowed through thoughtfully engineered software. The foundations established in this specification and design document provide a rigorous, extensible baseline for the full production implementation and multi-institution deployment planned for the 2026--2027 academic year.
		\end{tcolorbox}
		
	\end{spacing}
	
\end{document}